\section{Exactness and Error Propagation}
\label{sec:theory}

The analysis treats the supplied gate as an arbitrary fixed relation.
It separates compiler correctness from the semantic accuracy of its inputs.
Complete proofs and boundary cases appear in Appendix~\ref{app:proofs}.

\begin{theorem}[All-cutoff equivalence]
\label{thm:cutoff}
For every dated claim $c$ and finite cutoff $\tau$,
\begin{equation}
c\in A_G(\tau)\quad\Longleftrightarrow\quad\tau\in I_G(c).
\label{eq:equivalence}
\end{equation}
At most one stored witness per claim therefore represents every temporal eligibility decision.
\end{theorem}
\begin{proof}
If $t_c>\tau$, both statements are false.
Otherwise, an accepted replacement exists by $\tau$ exactly when $b_G(c)\leq\tau$.
Taking the complement gives the equivalence.
\end{proof}

This result also explains why a newest replacement is insufficient for historical retrieval.
A newest witness establishes eventual replacement, but can miss an earlier endpoint.
Both choices nevertheless give the same final current mask.
Each excludes exactly the claims with at least one accepted later witness.

\begin{corollary}[Temporal prefix consistency]
\label{cor:prefix}
Let $\mathcal C_{\leq\tau}=\{c:t_c\leq\tau\}$.
Querying full-corpus certificates at $\tau$ equals compiling the final mask on $\mathcal C_{\leq\tau}$.
\end{corollary}
\begin{proof}
Both procedures examine precisely the accepted pairs satisfying $t_c<t_d\leq\tau$.
\end{proof}

Appending records dated after $\tau$ therefore preserves the evidence state at $\tau$.
This statement concerns timestamp order, not file arrival order.
A newly received record with an earlier timestamp can legitimately change the corresponding historical state.
The gate must also remain fixed when the corpus grows.

\begin{theorem}[Pair-local influence]
\label{thm:influence}
Fix claim timestamps and identifiers.
Suppose gates $G$ and $H$ differ on $h$ temporally ordered pairs.
Let $U$ contain the older endpoints of those pairs.
For every cutoff,
\begin{equation}
A_G(\tau)\triangle A_H(\tau)\subseteq U,
\qquad |U|\leq h.
\label{eq:influence}
\end{equation}
With a fixed ranking and deterministic ties,
\begin{equation}
|R_G(q,\tau)\triangle R_H(q,\tau)|
\leq 2\min\{k,|U|\}.
\label{eq:context-bound}
\end{equation}
\end{theorem}
\begin{proof}
A claim's boundary depends only on pairs whose older endpoint is that claim.
Its interval cannot change outside $U$.
Changing one eligibility bit removes or adds at most one selected claim and one replacement.
Apply this argument successively over $U$.
The two contexts also contain at most $2k$ claims altogether.
\end{proof}

For one changed pair, the interval difference is confined to one record.
Its temporal extent is also explicit.
Writing $J_c=[\min\{b_G(c),b_H(c)\},\max\{b_G(c),b_H(c)\})$, we have
\begin{equation}
c\in A_G(\tau)\triangle A_H(\tau)
\quad\Longleftrightarrow\quad\tau\in J_c.
\label{eq:footprint}
\end{equation}
For a query-time distribution $\mu$, the expected mask difference equals $\sum_c\mu(J_c)$.
This identifies the period affected by each changed decision.
The bound concerns pair decisions, so a parser update can affect several intervals through several changed pairs.

\paragraph{Why component closure differs.}
An incorrect bridge can change intervals throughout a connected component.
Consider one attribute with repeated value $a$ at times $0,\ldots,m-1$.
It changes to $b$ at time $2m$.
A second attribute has $z$ at time $m$ and $w$ at time $2m+1$.
All four values differ.
An incorrect accepted pair between $b$ and $w$ merges the histories.
Component sorting then closes every earlier $a$ record at time $m$.
At cutoff $m+\tfrac12$, it excludes $m$ records through one changed pair.
Direct compilation changes none at that cutoff because the affected older endpoint has not appeared.
This construction gives unbounded component amplification and a failure of temporal prefix consistency.
It is a correctness example, not a benchmark observation.

\begin{proposition}[Compiler exactness]
\label{prop:compiler}
The batched compiler returns Equation~\ref{eq:boundary} when its candidate postings contain every pair accepted by the gate.
\end{proposition}
\begin{proof}
Before each batch, the index contains exactly the earlier claims without an accepted earlier witness.
Completeness exposes every accepted pair from the batch to those claims.
The first such batch supplies the minimum replacement time.
Removing resolved claims cannot change that minimum.
Inserting after comparison prevents equal-time replacement.
\end{proof}
